\documentclass[aspectratio=169,11pt]{beamer}

% Compile with XeLaTeX
\usepackage{fontspec}
\setsansfont{DejaVu Sans}
\setmonofont{DejaVu Sans Mono}
\usepackage{polyglossia}
\setdefaultlanguage{english}
\usepackage{booktabs}
\usepackage{array}
\usepackage{tabularx}
\usepackage{adjustbox}
\usepackage{listings}
\usepackage{xcolor}
\usepackage{ragged2e}
\usepackage{graphicx}
\usepackage{hyperref}

\graphicspath{{images/}{./}{/mnt/data/images/}}
\hypersetup{unicode=true}

\usetheme{Madrid}
\usecolortheme{seahorse}
\setbeamertemplate{navigation symbols}{}
\setbeamertemplate{footline}[frame number]
\setbeamertemplate{itemize items}[circle]
\setbeamertemplate{enumerate items}[default]

\definecolor{codebg}{RGB}{248,248,248}
\definecolor{codekw}{RGB}{0,72,160}
\definecolor{codestr}{RGB}{150,60,30}
\definecolor{accent}{RGB}{0,90,150}
\definecolor{placeholderbg}{RGB}{245,248,252}
\definecolor{placeholderline}{RGB}{90,120,150}

\lstset{
  basicstyle=\ttfamily\scriptsize,
  keywordstyle=\color{codekw}\bfseries,
  stringstyle=\color{codestr},
  backgroundcolor=\color{codebg},
  frame=single,
  breaklines=true,
  columns=fullflexible,
  showstringspaces=false,
  keepspaces=true,
  language=SQL
}

\newcommand{\keyword}[1]{\textcolor{accent}{\textbf{#1}}}
\newcommand{\kw}[1]{\keyword{#1}}
\newcommand{\tightlist}{\setlength{\itemsep}{2pt}\setlength{\parskip}{0pt}}
\newcommand{\smallitem}{\setlength{\itemsep}{2pt}\setlength{\parskip}{0pt}}
\newcommand{\qoption}[2]{\item[#1.] #2}

% Image placeholder: if the image exists, it is inserted; otherwise a box is shown.
% To insert real images later, place PNG files into the images/ folder with the exact filenames.
\newcommand{\imageplaceholder}[2]{%
\IfFileExists{#1}{\includegraphics[width=0.78\textwidth,height=0.58\textheight,keepaspectratio]{#1}}{%
\IfFileExists{images/#1}{\includegraphics[width=0.78\textwidth,height=0.58\textheight,keepaspectratio]{images/#1}}{%
\fcolorbox{placeholderline}{placeholderbg}{\parbox[c][0.42\textheight][c]{0.78\textwidth}{\centering\Large Illustration not found\\[2mm]\normalsize #2\\[2mm]\scriptsize Check file: \texttt{images/\detokenize{#1}}}}}}}

\title[ER Model in DBMS]{Introduction to the ER Model in DBMS}
\subtitle{Entity, Attribute, Relationship, Degree, Cardinality, and Participation Constraint}
\author{Lecture Material}
\institute{Course: Databases}
\date{Updated: 02/04/2026}

\AtBeginSection[]{
  \begin{frame}
    \centering
    \vfill
    {\Huge \insertsection}\par
    \vfill
  \end{frame}
}

\begin{document}

%------------------------------------------------
\begin{frame}
  \titlepage
\end{frame}

%------------------------------------------------
\begin{frame}{Learning Objectives}
\small
After completing this lecture, learners will be able to:
\begin{enumerate}
\smallitem
  \item Explain the concept of the \kw{Entity-Relationship Model} in database design.
  \item Describe the role of the \kw{ER Model} and \kw{ER Diagram}.
  \item Distinguish the main components: \kw{entity}, \kw{attribute}, and \kw{relationship}.
  \item Identify \kw{strong entities} and \kw{weak entities}.
  \item Distinguish attribute types: simple, key, composite, multivalued, and derived.
  \item Explain \kw{degree}, \kw{cardinality}, and \kw{participation constraint}.
  \item Apply the basic process for drawing an ER Diagram for a practical problem.
\end{enumerate}
\end{frame}

%------------------------------------------------
\begin{frame}{Lecture Outline}
\small
\begin{columns}[T]
\begin{column}{0.48\textwidth}
\begin{enumerate}
\smallitem
  \item Overview of the ER Model
  \item Role in database design
  \item Components of the ER Model
  \item Entity, strong entity, and weak entity
  \item Attribute and attribute types
\end{enumerate}
\end{column}
\begin{column}{0.48\textwidth}
\begin{enumerate}
\setcounter{enumi}{5}
\smallitem
  \item Relationship and degree
  \item Cardinality
  \item Participation constraint
  \item How to draw an ER Diagram
  \item Examples, questions, and exercises
\end{enumerate}
\end{column}
\end{columns}
\end{frame}

%================================================
\section{Overview of the ER Model}

%------------------------------------------------
\begin{frame}{What is the ER Model?}
\small
The \kw{Entity-Relationship Model}, or \kw{ER Model}, is a conceptual model used to design databases.

\vspace{2mm}
This model represents the logical structure of a database through:
\begin{itemize}
\smallitem
  \item The \kw{entities} whose data must be stored.
  \item The \kw{attributes} that describe those entities.
  \item The \kw{relationships} among entities.
\end{itemize}
\end{frame}

%------------------------------------------------
\begin{frame}{What Questions Does the ER Model Answer?}
\small
The ER Model helps designers answer:
\begin{itemize}
\smallitem
  \item What objects does the system need to store?
  \item What information does each object have?
  \item How are the objects related to one another?
  \item What constraints should be represented in the data model?
\end{itemize}

\begin{block}{ER Diagram}
The graphical representation of the ER Model is called an \kw{Entity-Relationship Diagram}, abbreviated as \kw{ERD}.
\end{block}
\end{frame}

%------------------------------------------------
\begin{frame}{Illustration: Overview of the ER Model}
\centering
\imageplaceholder{introduction-to-er-model.png}{Overview of the Entity-Relationship Model}

\vspace{2mm}
{\small Figure 1. Overview of the ER Model.}
\end{frame}

%------------------------------------------------
\begin{frame}{Role of the ER Model in Database Design}
\small
The database design process usually involves three stages:
\begin{enumerate}
\smallitem
  \item \kw{Requirements gathering}: asking users, customers, or stakeholders about system functions and data requirements.
  \item \kw{Conceptual or logical design}: using the ER Model to describe entities, attributes, and relationships.
  \item \kw{Physical design}: defining tables, keys, indexes, views, and implementation details for a specific DBMS.
\end{enumerate}
\end{frame}

%------------------------------------------------
\begin{frame}{Why Are ER Diagrams Useful?}
\small
ER Diagrams are useful because they:
\begin{itemize}
\smallitem
  \item Represent data and data relationships visually.
  \item Can be converted into relational tables.
  \item Model real-world objects.
  \item Do not require viewers to deeply understand a specific DBMS.
  \item Make complex systems easier to understand.
\end{itemize}
\end{frame}

%------------------------------------------------
\begin{frame}{Quick Quiz 1: Overview of the ER Model}
\small
\textbf{Question 1.} What is the ER Model mainly used for?
\begin{enumerate}
\tightlist
  \qoption{A}{Designing the conceptual model of a database}
  \qoption{B}{Increasing the CPU speed of a server}
  \qoption{C}{Designing user interfaces}
  \qoption{D}{Compressing image data}
\end{enumerate}

\textbf{Question 2.} What is an ER Diagram?
\begin{enumerate}
\tightlist
  \qoption{A}{A data backup tool}
  \qoption{B}{A graphical representation of the ER Model}
  \qoption{C}{A database operating system}
  \qoption{D}{A web programming language}
\end{enumerate}
\end{frame}

%------------------------------------------------
\begin{frame}{Quick Quiz 1: Overview of the ER Model (cont.)}
\small
\textbf{Question 3.} In which design step is the ER Model most commonly used?
\begin{enumerate}
\tightlist
  \qoption{A}{Network card configuration}
  \qoption{B}{Source code compilation}
  \qoption{C}{Conceptual or logical design}
  \qoption{D}{Deleting old data}
\end{enumerate}
\end{frame}

%================================================
\section{Main Components of the ER Model}

%------------------------------------------------
\begin{frame}{Three Core Components}
\small
The ER Model consists of three core components:
\vspace{2mm}

\begin{tabularx}{\textwidth}{>{\bfseries}p{0.22\textwidth} X p{0.28\textwidth}}
\toprule
Component & Meaning & Common ERD Symbol \\
\midrule
Entity & An object or concept whose data must be stored & Rectangle \\
Attribute & A property that describes an entity & Oval \\
Relationship & A connection between entities & Diamond \\
\bottomrule
\end{tabularx}
\end{frame}

%------------------------------------------------
\begin{frame}{Illustration: Components of an ER Diagram}
\centering
\imageplaceholder{components-of-er-diagram.png}{Basic components of an ER Diagram}

\vspace{2mm}
{\small Figure 2. Entity, Attribute, and Relationship in an ER Diagram.}
\end{frame}

%================================================
\section{Entity in the ER Model}

%------------------------------------------------
\begin{frame}{What is an Entity?}
\small
An \kw{Entity} is an object, thing, concept, or event in the real world about which the system needs to store information.

\vspace{2mm}
Examples:
\begin{columns}[T]
\begin{column}{0.48\textwidth}
\begin{itemize}
\smallitem
  \item \texttt{Student}
  \item \texttt{Course}
  \item \texttt{Employee}
  \item \texttt{Company}
\end{itemize}
\end{column}
\begin{column}{0.48\textwidth}
\begin{itemize}
\smallitem
  \item \texttt{Product}
  \item \texttt{Reservation}
  \item \texttt{Device}
  \item \texttt{Document}
\end{itemize}
\end{column}
\end{columns}

\vspace{2mm}
In a relational database, an entity is usually converted into a \kw{table}.
\end{frame}

%------------------------------------------------
\begin{frame}{Entity Type, Entity Instance, and Entity Set}
\small
\begin{block}{Entity Type}
Describes the general structure of one type of entity, for example: \texttt{Student(student\_id, full\_name, date\_of\_birth, email)}.
\end{block}

\begin{block}{Entity Instance}
A specific object belonging to an entity type, for example: \texttt{S001, Nguyen An, 2005-01-15, an@example.com}.
\end{block}

\begin{block}{Entity Set}
The collection of all entities of the same entity type, for example all students in a university.
\end{block}
\end{frame}

%------------------------------------------------
\begin{frame}{Reminder about Entity Sets}
\small
An ER Diagram usually represents an \kw{entity type} or an \kw{entity set}, not each specific data row.

\vspace{2mm}
\begin{exampleblock}{Example}
An ERD represents the structure \texttt{Student}; it does not draw each specific student such as S001, S002, or S003.
\end{exampleblock}
\end{frame}

%------------------------------------------------
\begin{frame}{Illustration: Entity Set}
\centering
\imageplaceholder{entity-set.png}{Entity type, entity instance, and entity set}

\vspace{2mm}
{\small Figure 3. Entity type, entity instance, and entity set.}
\end{frame}

%------------------------------------------------
\begin{frame}{Strong Entity}
\small
A \kw{Strong Entity} is an entity type that can be uniquely identified by its own attributes.

\vspace{2mm}
Characteristics:
\begin{itemize}
\smallitem
  \item It has a \kw{key attribute} or \kw{primary key}.
  \item It does not depend on another entity for identification.
  \item It is usually represented by a \kw{single rectangle} in an ER Diagram.
\end{itemize}

\begin{exampleblock}{Example}
\texttt{Student} has \texttt{student\_id}; \texttt{Employee} has \texttt{employee\_id}; \texttt{Product} has \texttt{product\_id}.
\end{exampleblock}
\end{frame}

%------------------------------------------------
\begin{frame}{Weak Entity}
\small
A \kw{Weak Entity} is an entity type that cannot be uniquely identified by its own attributes alone.

\vspace{2mm}
A Weak Entity must depend on a \kw{Strong Entity} for identification.
\begin{itemize}
\smallitem
  \item It does not have a sufficiently strong key of its own.
  \item It depends on an identifying strong entity.
  \item It usually has \kw{total participation} in its relationship with the strong entity.
  \item It is represented by a \kw{double rectangle}.
  \item Its identifying relationship is usually represented by a \kw{double diamond}.
\end{itemize}
\end{frame}

%------------------------------------------------
\begin{frame}{Example of a Weak Entity}
\small
A company stores information about employees' dependents.

\vspace{2mm}
\begin{itemize}
\smallitem
  \item \texttt{Employee} is a strong entity.
  \item \texttt{Dependent} is a weak entity.
  \item \texttt{Dependent} cannot exist independently without being linked to a specific employee.
\end{itemize}
\end{frame}

%------------------------------------------------
\begin{frame}{Illustration: Strong Entity and Weak Entity}
\centering
\imageplaceholder{strong-weak-entity.png}{Strong Entity and Weak Entity}

\vspace{2mm}
{\small Figure 4. Strong Entity and Weak Entity.}
\end{frame}

%------------------------------------------------
\begin{frame}{Quick Quiz 2: Entity}
\small
\textbf{Question 4.} What is an entity in the ER Model?
\begin{enumerate}
\tightlist
  \qoption{A}{A numeric data type}
  \qoption{B}{An index type}
  \qoption{C}{An SQL statement}
  \qoption{D}{A real-world object or concept whose data must be stored}
\end{enumerate}

\textbf{Question 5.} Which of the following is usually an example of a strong entity?
\begin{enumerate}
\tightlist
  \qoption{A}{A secondary phone number with no owner}
  \qoption{B}{A temporary log row with no identifier}
  \qoption{C}{\texttt{Student} with a unique \texttt{student\_id}}
  \qoption{D}{An attribute derived from date of birth}
\end{enumerate}
\end{frame}

%------------------------------------------------
\begin{frame}{Quick Quiz 2: Entity (cont.)}
\small
\textbf{Question 6.} Which characteristic belongs to a Weak Entity?
\begin{enumerate}
\tightlist
  \qoption{A}{It is always a multivalued attribute}
  \qoption{B}{It depends on a strong entity for identification}
  \qoption{C}{It cannot appear in an ER Diagram}
  \qoption{D}{It is always represented by a single oval}
\end{enumerate}
\end{frame}

%================================================
\section{Attribute in the ER Model}

%------------------------------------------------
\begin{frame}{What is an Attribute?}
\small
An \kw{Attribute} is a property used to describe an entity.

\vspace{2mm}
For example, the entity \texttt{Student} may have attributes such as:
\begin{itemize}
\smallitem
  \item \texttt{Roll\_No}
  \item \texttt{Name}
  \item \texttt{DOB}
  \item \texttt{Age}
  \item \texttt{Address}
  \item \texttt{Mobile\_No}
\end{itemize}

In an ER Diagram, an attribute is usually represented by an \kw{oval}.
\end{frame}

%------------------------------------------------
\begin{frame}{Illustration: Attribute}
\centering
\imageplaceholder{attribute.png}{Attribute in the ER Model}
\end{frame}

%------------------------------------------------
\begin{frame}{Types of Attributes}
\small
Common attribute types include:
\begin{enumerate}
\smallitem
  \item \kw{Simple Attribute}
  \item \kw{Key Attribute}
  \item \kw{Composite Attribute}
  \item \kw{Multivalued Attribute}
  \item \kw{Derived Attribute}
\end{enumerate}
\end{frame}

%------------------------------------------------
\begin{frame}{Simple Attribute}
\small
A \kw{Simple Attribute} cannot be divided into smaller attributes with independent meaning.

\vspace{2mm}
Examples:
\begin{itemize}
\smallitem
  \item \texttt{Age}
  \item \texttt{Gender}
  \item \texttt{Roll\_No}
\end{itemize}

\vspace{2mm}
\begin{alertblock}{Note about the illustration}
The original reference does not provide a separate illustration for Simple Attribute. In an ER Diagram, a simple attribute is still represented by a single oval, similar to ordinary attributes such as \texttt{Name}, \texttt{DOB}, or \texttt{Age}.
\end{alertblock}
\end{frame}

%------------------------------------------------
\begin{frame}{Key Attribute}
\small
A \kw{Key Attribute} is used to uniquely identify each entity in an entity set.

\vspace{2mm}
Examples:
\begin{itemize}
\smallitem
  \item \texttt{Roll\_No} uniquely identifies each student.
  \item \texttt{Employee\_ID} uniquely identifies each employee.
  \item \texttt{Product\_ID} uniquely identifies each product.
\end{itemize}

In an ER Diagram, a key attribute is represented by an underlined oval.
\end{frame}

%------------------------------------------------
\begin{frame}{Illustration: Key Attribute}
\centering
\imageplaceholder{key-attribute.png}{Key Attribute in the ER Model}
\end{frame}

%------------------------------------------------
\begin{frame}{Composite Attribute}
\small
A \kw{Composite Attribute} can be divided into smaller attributes.

\vspace{2mm}
For example, \texttt{Address} may consist of:
\begin{itemize}
\smallitem
  \item \texttt{Street}
  \item \texttt{City}
  \item \texttt{State}
  \item \texttt{Country}
\end{itemize}
\end{frame}

%------------------------------------------------
\begin{frame}{Illustration: Composite Attribute}
\centering
\imageplaceholder{simple-attribute.png}{Composite Attribute in the ER Model}
\end{frame}

%------------------------------------------------
\begin{frame}{Multivalued Attribute}
\small
A \kw{Multivalued Attribute} can have more than one value for a single entity.

\vspace{2mm}
Examples:
\begin{itemize}
\smallitem
  \item A student may have several phone numbers.
  \item An employee may have several skills.
  \item A product may have several colors.
\end{itemize}

In an ER Diagram, a multivalued attribute is represented by a double oval.
\end{frame}

%------------------------------------------------
\begin{frame}{Illustration: Multivalued Attribute}
\centering
\imageplaceholder{composite-attribute.png}{Multivalued Attribute in the ER Model}
\end{frame}

%------------------------------------------------
\begin{frame}{Derived Attribute}
\small
A \kw{Derived Attribute} can be derived from another attribute.

\vspace{2mm}
Examples:
\begin{itemize}
\smallitem
  \item \texttt{Age} can be derived from \texttt{DOB}.
  \item \texttt{Total\_Amount} can be calculated from quantity and unit price.
  \item \texttt{Years\_of\_Service} can be calculated from the starting date.
\end{itemize}

In an ER Diagram, a derived attribute is represented by a dashed oval.
\end{frame}

%------------------------------------------------
\begin{frame}{Illustration: Derived Attribute}
\centering
\imageplaceholder{multivalued-attribute.png}{Derived Attribute in the ER Model}
\end{frame}

%------------------------------------------------
\begin{frame}{Entity and Combined Attributes}
\small
When attribute types are combined in one ER Diagram, the entity \texttt{Student} may have:
\begin{itemize}
\smallitem
  \item \texttt{Roll\_No} as a key attribute.
  \item \texttt{Name} and \texttt{DOB} as simple attributes.
  \item \texttt{Phone\_No} as a multivalued attribute.
  \item \texttt{Age} as a derived attribute.
  \item \texttt{Address} as a composite attribute consisting of \texttt{Street}, \texttt{City}, \texttt{State}, and \texttt{Country}.
\end{itemize}
\end{frame}

%------------------------------------------------
\begin{frame}{Illustration: Entity and Combined Attributes}
\centering
\imageplaceholder{derived-attribute.png}{Student entity and attribute types in the ER Model}
\end{frame}

%------------------------------------------------
\begin{frame}{Quick Quiz 3: Attribute}
\small
\textbf{Question 7.} What is an attribute used for in the ER Model?
\begin{enumerate}
\tightlist
  \qoption{A}{Deleting the whole database}
  \qoption{B}{Increasing disk capacity}
  \qoption{C}{Describing a property of an entity}
  \qoption{D}{Replacing the operating system}
\end{enumerate}

\textbf{Question 8.} \texttt{Roll\_No} uniquely identifies a student. What type of attribute is it?
\begin{enumerate}
\tightlist
  \qoption{A}{Composite Attribute}
  \qoption{B}{Multivalued Attribute}
  \qoption{C}{Derived Attribute}
  \qoption{D}{Key Attribute}
\end{enumerate}
\end{frame}

%------------------------------------------------
\begin{frame}{Quick Quiz 3: Attribute (cont.)}
\small
\textbf{Question 9.} \texttt{Address} consisting of \texttt{Street}, \texttt{City}, \texttt{State}, and \texttt{Country} is an example of which attribute type?
\begin{enumerate}
\tightlist
  \qoption{A}{Composite Attribute}
  \qoption{B}{Derived Attribute}
  \qoption{C}{Multivalued Attribute}
  \qoption{D}{Weak Attribute}
\end{enumerate}

\textbf{Question 10.} \texttt{Age} calculated from \texttt{DOB} is an example of which attribute type?
\begin{enumerate}
\tightlist
  \qoption{A}{Key Attribute}
  \qoption{B}{Derived Attribute}
  \qoption{C}{Composite Attribute}
  \qoption{D}{Multivalued Attribute}
\end{enumerate}
\end{frame}

%================================================
\section{Relationship and Degree}

%------------------------------------------------
\begin{frame}{Relationship in the ER Model}
\small
A \kw{Relationship} describes a connection between entities.

\vspace{2mm}
Examples:
\begin{itemize}
\smallitem
  \item \texttt{Student} \kw{enrolls in} \texttt{Course}.
  \item \texttt{Employee} \kw{works for} \texttt{Department}.
  \item \texttt{Customer} \kw{places} \texttt{Order}.
  \item \texttt{Doctor} \kw{performs} \texttt{Surgery}.
\end{itemize}

In an ER Diagram, a relationship is usually represented by a diamond and connected to entities by lines.
\end{frame}

%------------------------------------------------
\begin{frame}{Relationship Type and Relationship Set}
\small
\begin{block}{Relationship Type}
Describes the type of association between entity types. Example: \texttt{Student -- Enrolled in -- Course}.
\end{block}

\begin{block}{Relationship Set}
The set of relationships of the same type, for example S1 enrolled in C2, S2 enrolled in C1, and S3 enrolled in C3.
\end{block}
\end{frame}

%------------------------------------------------
\begin{frame}{Illustration: Relationship Type}
\centering
\imageplaceholder{relationship-set.png}{Relationship Type between Student and Course}

\vspace{2mm}
{\small Figure 5a. Relationship type between Student and Course.}
\end{frame}

%------------------------------------------------
\begin{frame}{Illustration: Relationship Set}
\centering
\imageplaceholder{recursive-relationship.png}{Relationship set of the Enrolled in type}

\vspace{2mm}
{\small Figure 5b. Relationship set of the Enrolled in type.}
\end{frame}

%------------------------------------------------
\begin{frame}{Degree of a Relationship Set}
\small
The \kw{Degree} of a relationship set is the number of different entity sets participating in the relationship set.

\vspace{2mm}
There are four common forms:
\begin{enumerate}
\smallitem
  \item \kw{Unary / Recursive Relationship}: one participating entity set.
  \item \kw{Binary Relationship}: two participating entity sets.
  \item \kw{Ternary Relationship}: three participating entity sets.
  \item \kw{N-ary Relationship}: n participating entity sets.
\end{enumerate}
\end{frame}

%------------------------------------------------
\begin{frame}{Examples of Degree}
\small
\begin{tabularx}{\textwidth}{>{\bfseries}p{0.28\textwidth} X}
\toprule
Degree Type & Example \\
\midrule
Unary / Recursive & One employee manages another employee \\
Binary & A student enrolls in a course \\
Ternary & A supplier supplies a part to a project \\
N-ary & A relationship involving more than three entity sets \\
\bottomrule
\end{tabularx}
\end{frame}

%------------------------------------------------
\begin{frame}{Illustration: Unary / Recursive Relationship}
\centering
\imageplaceholder{binary-relationship.png}{Unary or Recursive Relationship}
\end{frame}

%------------------------------------------------
\begin{frame}{Illustration: Binary Relationship}
\centering
\imageplaceholder{relationship-set.png}{Binary Relationship}
\end{frame}

%------------------------------------------------
\begin{frame}{Illustration: Ternary Relationship}
\centering
\imageplaceholder{ternary-relationship.png}{Ternary Relationship}
\end{frame}

%------------------------------------------------
\begin{frame}{Illustration: N-ary Relationship}
\centering
\imageplaceholder{n-ary-relationship.png}{N-ary Relationship}
\end{frame}

%------------------------------------------------
\begin{frame}{Quick Quiz 4: Relationship and Degree}
\small
\textbf{Question 11.} What does a relationship represent in the ER Model?
\begin{enumerate}
\tightlist
  \qoption{A}{A font type}
  \qoption{B}{A connection between entities}
  \qoption{C}{A CPU performance metric}
  \qoption{D}{An image data type}
\end{enumerate}

\textbf{Question 12.} What is a relationship involving two entity sets called?
\begin{enumerate}
\tightlist
  \qoption{A}{Unary Relationship}
  \qoption{B}{Ternary Relationship}
  \qoption{C}{Binary Relationship}
  \qoption{D}{N-ary Relationship}
\end{enumerate}
\end{frame}

%------------------------------------------------
\begin{frame}{Quick Quiz 4: Relationship and Degree (cont.)}
\small
\textbf{Question 13.} A relationship in which one employee manages another employee is an example of which degree?
\begin{enumerate}
\tightlist
  \qoption{A}{Ternary Relationship}
  \qoption{B}{Binary Relationship}
  \qoption{C}{N-ary Relationship}
  \qoption{D}{Unary / Recursive Relationship}
\end{enumerate}

\textbf{Question 14.} What does the degree of a relationship set indicate?
\begin{enumerate}
\tightlist
  \qoption{A}{The number of entity sets participating in the relationship set}
  \qoption{B}{The number of records in the database}
  \qoption{C}{The number of columns in a table}
  \qoption{D}{The number of indexes created}
\end{enumerate}
\end{frame}

%================================================
\section{Cardinality in the ER Model}

%------------------------------------------------
\begin{frame}{What is Cardinality?}
\small
\kw{Cardinality} indicates the maximum number of times an entity in an entity set can participate in a relationship set.

\vspace{2mm}
Cardinality helps answer:
\begin{itemize}
\smallitem
  \item How many other entities can one entity be associated with?
  \item Is the relationship one-to-one, one-to-many, or many-to-many?
  \item When converting to relational tables, do we need foreign keys or an intermediate table?
\end{itemize}
\end{frame}

%------------------------------------------------
\begin{frame}{One-to-One Relationship}
\small
A \kw{one-to-one} relationship occurs when each entity in each entity set participates at most once in the relationship.

\vspace{2mm}
Examples:
\begin{itemize}
\smallitem
  \item One person has only one passport.
  \item One passport belongs to only one person.
\end{itemize}

\begin{center}
\texttt{Person 1 -- 1 Passport}
\end{center}
\end{frame}

%------------------------------------------------
\begin{frame}{Illustration: One-to-One}
\centering
\imageplaceholder{one-to-one.png}{One-to-One Relationship}
\end{frame}

%------------------------------------------------
\begin{frame}{Set Representation: One-to-One}
\centering
\imageplaceholder{one-to-one-set-representation.png}{Set Representation of a One-to-One Relationship}
\end{frame}

%------------------------------------------------
\begin{frame}{One-to-Many Relationship}
\small
A \kw{one-to-many} relationship occurs when one entity on the first side can be associated with multiple entities on the second side.

\vspace{2mm}
Examples:
\begin{itemize}
\smallitem
  \item One department has many doctors.
  \item One customer has many orders.
  \item One class has many students.
\end{itemize}

\begin{center}
\texttt{Department 1 -- M Doctor}
\end{center}
\end{frame}

%------------------------------------------------
\begin{frame}{Illustration: One-to-Many}
\centering
\imageplaceholder{one-to-many.png}{One-to-Many Relationship}
\end{frame}

%------------------------------------------------
\begin{frame}{Set Representation: One-to-Many}
\centering
\imageplaceholder{one-to-many-set-representation.png}{Set Representation of a One-to-Many Relationship}
\end{frame}

%------------------------------------------------
\begin{frame}{Many-to-One Relationship}
\small
A \kw{many-to-one} relationship is the reverse direction of one-to-many.

\vspace{2mm}
Examples:
\begin{itemize}
\smallitem
  \item Many surgeries can be performed by one surgeon.
  \item Many orders belong to one customer.
  \item Many employees belong to one department.
\end{itemize}

\begin{center}
\texttt{Surgery M -- 1 Surgeon}
\end{center}
\end{frame}

%------------------------------------------------
\begin{frame}{Illustration: Many-to-One}
\centering
\imageplaceholder{many-to-one.png}{Many-to-One Relationship}
\end{frame}

%------------------------------------------------
\begin{frame}{Set Representation: Many-to-One}
\centering
\imageplaceholder{many-to-one-set-representation.png}{Set Representation of a Many-to-One Relationship}
\end{frame}

%------------------------------------------------
\begin{frame}{Many-to-Many Relationship}
\small
A \kw{many-to-many} relationship occurs when entities on both sides can participate multiple times in the relationship.

\vspace{2mm}
Examples:
\begin{itemize}
\smallitem
  \item A student studies many courses.
  \item A course has many students.
  \item An employee works on many projects.
  \item A project has many employees.
\end{itemize}
\end{frame}

%------------------------------------------------
\begin{frame}[fragile]{Intermediate Table for Many-to-Many}
\small
When converting to the relational model, a many-to-many relationship usually requires an intermediate table.

\begin{verbatim}
Students(student_id, full_name)
Courses(course_id, course_name)
Enrollments(student_id, course_id, semester, grade)
\end{verbatim}
\end{frame}

%------------------------------------------------
\begin{frame}{Illustration: Many-to-Many}
\centering
\imageplaceholder{many-to-many.png}{Many-to-Many Relationship}
\end{frame}

%------------------------------------------------
\begin{frame}{Set Representation: Many-to-Many}
\centering
\imageplaceholder{many-to-many-set-representation.png}{Set Representation of a Many-to-Many Relationship}
\end{frame}

%------------------------------------------------
\begin{frame}{Quick Quiz 5: Cardinality}
\small
\textbf{Question 15.} What does cardinality indicate in the ER Model?
\begin{enumerate}
\tightlist
  \qoption{A}{The maximum number of times an entity can participate in a relationship set}
  \qoption{B}{The color of a rectangle in an ERD}
  \qoption{C}{The number of backup files}
  \qoption{D}{The SQL query speed}
\end{enumerate}

\textbf{Question 16.} One person has one passport, and one passport belongs to one person. What is this relationship?
\begin{enumerate}
\tightlist
  \qoption{A}{One-to-Many}
  \qoption{B}{Many-to-One}
  \qoption{C}{One-to-One}
  \qoption{D}{Many-to-Many}
\end{enumerate}
\end{frame}

%------------------------------------------------
\begin{frame}{Quick Quiz 5: Cardinality (cont.)}
\small
\textbf{Question 17.} One customer has many orders. What relationship type is this?
\begin{enumerate}
\tightlist
  \qoption{A}{One-to-One}
  \qoption{B}{Many-to-Many}
  \qoption{C}{Many-to-One}
  \qoption{D}{One-to-Many}
\end{enumerate}

\textbf{Question 18.} If a student can take many courses and a course can have many students, what is the usual relationship type?
\begin{enumerate}
\tightlist
  \qoption{A}{One-to-One}
  \qoption{B}{Many-to-Many}
  \qoption{C}{One-to-Many}
  \qoption{D}{Recursive}
\end{enumerate}
\end{frame}

%================================================
\section{Participation Constraint}

%------------------------------------------------
\begin{frame}{What is Participation Constraint?}
\small
A \kw{Participation Constraint} describes whether an entity in an entity set is required to participate in a relationship.

\vspace{2mm}
There are two main types:
\begin{enumerate}
\smallitem
  \item \kw{Total Participation}
  \item \kw{Partial Participation}
\end{enumerate}
\end{frame}

%------------------------------------------------
\begin{frame}{Total Participation}
\small
\kw{Total Participation} means that every entity in the entity set must participate in the relationship.

\vspace{2mm}
Example: if the system requires every student to enroll in at least one course, then the entity set \texttt{Student} has total participation in the relationship \texttt{Enrolled in}.

\vspace{2mm}
In an ER Diagram, total participation is usually represented by a \kw{double line}.
\end{frame}

%------------------------------------------------
\begin{frame}{Illustration: Total Participation}
\centering
\imageplaceholder{total-participation.png}{Total Participation in the ER Model}
\end{frame}

%------------------------------------------------
\begin{frame}{Partial Participation}
\small
\kw{Partial Participation} means that an entity in the entity set may or may not participate in the relationship.

\vspace{2mm}
Example: some courses may not have any enrolled students yet. In that case, the entity set \texttt{Course} has partial participation in the relationship \texttt{Enrolled in}.

\vspace{2mm}
In an ER Diagram, partial participation is usually represented by a \kw{single line}.
\end{frame}

%------------------------------------------------
\begin{frame}{Set Representation: Participation Constraint}
\centering
\imageplaceholder{participation-set-representation.png}{Set Representation of Participation Constraint}
\end{frame}

%------------------------------------------------
\begin{frame}{Quick Quiz 6: Participation Constraint}
\small
\textbf{Question 19.} What does Total Participation mean?
\begin{enumerate}
\tightlist
  \qoption{A}{An entity never participates in a relationship}
  \qoption{B}{Only one entity participates in a relationship}
  \qoption{C}{All entities in the entity set must participate in the relationship}
  \qoption{D}{The relationship is removed from the ERD}
\end{enumerate}

\textbf{Question 20.} What does Partial Participation mean?
\begin{enumerate}
\tightlist
  \qoption{A}{An entity may or may not participate in a relationship}
  \qoption{B}{An entity must participate in a relationship}
  \qoption{C}{An entity is always a weak entity}
  \qoption{D}{An entity always has many primary keys}
\end{enumerate}
\end{frame}

%================================================
\section{How to Draw an ER Diagram}

%------------------------------------------------
\begin{frame}{Process for Drawing an ER Diagram}
\small
An ER Diagram can be drawn in six steps:
\begin{enumerate}
\smallitem
  \item Identify entities.
  \item Identify relationships.
  \item Add attributes.
  \item Define primary keys.
  \item Remove redundancies.
  \item Review for clarity.
\end{enumerate}
\end{frame}

%------------------------------------------------
\begin{frame}{Step 1: Identify Entities}
\small
Identify all important entities in the system.

\vspace{2mm}
For example, in a student management system:
\begin{itemize}
\smallitem
  \item \texttt{Student}
  \item \texttt{Course}
  \item \texttt{Lecturer}
  \item \texttt{Class}
  \item \texttt{Department}
\end{itemize}

Represent entities using rectangles.
\end{frame}

%------------------------------------------------
\begin{frame}{Step 2: Identify Relationships}
\small
Identify relationships among entities.

\vspace{2mm}
Examples:
\begin{itemize}
\smallitem
  \item \texttt{Student} enrolls in \texttt{Course}.
  \item \texttt{Lecturer} teaches \texttt{Class}.
  \item \texttt{Department} manages \texttt{Lecturer}.
\end{itemize}

Represent relationships using diamonds.
\end{frame}

%------------------------------------------------
\begin{frame}{Step 3: Add Attributes}
\small
Attach attributes to entities using ovals.

\vspace{2mm}
For example, the entity \texttt{Student} may have:
\begin{itemize}
\smallitem
  \item \texttt{student\_id}
  \item \texttt{full\_name}
  \item \texttt{date\_of\_birth}
  \item \texttt{email}
  \item \texttt{phone\_number}
\end{itemize}
\end{frame}

%------------------------------------------------
\begin{frame}{Step 4: Define Primary Keys}
\small
For each entity, identify the attribute that uniquely identifies it.

\vspace{2mm}
Examples:
\begin{itemize}
\smallitem
  \item \texttt{Student}: \texttt{student\_id}
  \item \texttt{Course}: \texttt{course\_id}
  \item \texttt{Lecturer}: \texttt{lecturer\_id}
\end{itemize}

In an ER Diagram, key attributes are usually underlined.
\end{frame}

%------------------------------------------------
\begin{frame}{Steps 5 and 6: Finalize the ER Diagram}
\small
\begin{block}{Step 5: Remove Redundancies}
Remove unnecessary entities, duplicate relationships, repeated attributes, and relationships without business meaning.
\end{block}

\begin{block}{Step 6: Review for Clarity}
Check whether the diagram is easy to understand, can be converted into relational tables, and correctly reflects the system requirements.
\end{block}
\end{frame}

%================================================
\section{Worked Example: Course Registration System}

%------------------------------------------------
\begin{frame}{Example Problem}
\small
Consider a university course registration system.

\vspace{2mm}
We need to identify:
\begin{itemize}
\smallitem
  \item Entities.
  \item Attributes.
  \item Relationships.
  \item Cardinalities.
  \item How to convert the ER model into relational tables.
\end{itemize}
\end{frame}

%------------------------------------------------
\begin{frame}{Main Entities}
\small
Main entities:
\begin{itemize}
\smallitem
  \item \texttt{Student}
  \item \texttt{Course}
  \item \texttt{Lecturer}
  \item \texttt{Department}
\end{itemize}
\end{frame}

%------------------------------------------------
\begin{frame}{Attributes of the Entities}
\scriptsize
\begin{columns}[T]
\begin{column}{0.48\textwidth}
\textbf{Student}
\begin{itemize}
\smallitem
  \item \texttt{student\_id}
  \item \texttt{full\_name}
  \item \texttt{email}
  \item \texttt{date\_of\_birth}
\end{itemize}

\textbf{Course}
\begin{itemize}
\smallitem
  \item \texttt{course\_id}
  \item \texttt{course\_name}
  \item \texttt{credits}
\end{itemize}
\end{column}
\begin{column}{0.48\textwidth}
\textbf{Lecturer}
\begin{itemize}
\smallitem
  \item \texttt{lecturer\_id}
  \item \texttt{full\_name}
  \item \texttt{email}
\end{itemize}

\textbf{Department}
\begin{itemize}
\smallitem
  \item \texttt{department\_id}
  \item \texttt{department\_name}
\end{itemize}
\end{column}
\end{columns}
\end{frame}

%------------------------------------------------
\begin{frame}{Main Relationships}
\small
Main relationships:
\begin{itemize}
\smallitem
  \item \texttt{Student} \kw{enrolls in} \texttt{Course}.
  \item \texttt{Lecturer} \kw{teaches} \texttt{Course}.
  \item \texttt{Department} \kw{offers} \texttt{Course}.
  \item \texttt{Department} \kw{manages} \texttt{Lecturer}.
\end{itemize}
\end{frame}

%------------------------------------------------
\begin{frame}{Cardinality in the Example}
\small
Some cardinalities:
\begin{itemize}
\smallitem
  \item A student can enroll in many courses.
  \item A course can have many enrolled students.
  \item A lecturer can teach many courses.
  \item A course may be taught by one or more lecturers depending on the rules.
  \item A department can manage many lecturers.
\end{itemize}
\end{frame}

%------------------------------------------------
\begin{frame}[fragile]{Converting to Relational Tables}
\small
The ER model can be converted into the following tables:
\begin{verbatim}
Students(student_id, full_name, email, date_of_birth)
Courses(course_id, course_name, credits)
Lecturers(lecturer_id, full_name, email)
Departments(department_id, department_name)
Enrollments(student_id, course_id, semester, grade)
Teaches(lecturer_id, course_id, semester)
\end{verbatim}
\end{frame}

%------------------------------------------------
\begin{frame}{Interpreting Intermediate Tables}
\small
\begin{itemize}
\smallitem
  \item \texttt{student\_id} is the primary key of \texttt{Students}.
  \item \texttt{course\_id} is the primary key of \texttt{Courses}.
  \item \texttt{lecturer\_id} is the primary key of \texttt{Lecturers}.
  \item \texttt{department\_id} is the primary key of \texttt{Departments}.
  \item \texttt{Enrollments} is an intermediate table for the many-to-many relationship between \texttt{Student} and \texttt{Course}.
  \item \texttt{Teaches} links \texttt{Lecturer} and \texttt{Course}.
\end{itemize}
\end{frame}

%------------------------------------------------
\begin{frame}{Quick Quiz 7: Drawing an ERD}
\small
\textbf{Question 21.} When drawing an ER Diagram, what is usually the first step?
\begin{enumerate}
\tightlist
  \qoption{A}{Choosing the background color of the diagram}
  \qoption{B}{Creating indexes for tables}
  \qoption{C}{Deleting all relationships}
  \qoption{D}{Identifying the important entities}
\end{enumerate}

\textbf{Question 22.} In an ER Diagram, what symbol usually represents a relationship?
\begin{enumerate}
\tightlist
  \qoption{A}{Oval}
  \qoption{B}{Rectangle}
  \qoption{C}{Dashed line}
  \qoption{D}{Diamond}
\end{enumerate}
\end{frame}

%================================================
\section{Questions and Exercises}

%------------------------------------------------
\begin{frame}[allowframebreaks]{Short-Answer Review Questions}
\small
\begin{enumerate}
\smallitem
  \item What is the ER Model? Why is it considered a conceptual model in database design?
  \item Distinguish entity, attribute, and relationship.
  \item Distinguish strong entity and weak entity. Give an example.
  \item Present the common attribute types in the ER Model.
  \item What is the degree of a relationship set? Give examples of unary, binary, and ternary relationships.
  \item What is cardinality? Why is it important when converting an ER Model into a relational model?
  \item Distinguish total participation and partial participation.
\end{enumerate}
\end{frame}

%------------------------------------------------
\begin{frame}{Exercise 1: University System}
\small
A university needs to build a system for managing students, courses, lecturers, and departments.

\vspace{2mm}
\begin{enumerate}
\smallitem
  \item Identify the main entities.
  \item List some attributes for each entity.
  \item Identify relationships among the entities.
  \item Determine the cardinality for each relationship.
  \item Convert the ER model into a list of relational tables.
\end{enumerate}
\end{frame}

%------------------------------------------------
\begin{frame}{Exercise 2: Online Store}
\small
An online store needs to manage customers, products, orders, and payments.

\vspace{2mm}
\begin{enumerate}
\smallitem
  \item Identify the main entities and attributes.
  \item Identify relationships among \texttt{Customer}, \texttt{Order}, \texttt{Product}, and \texttt{Payment}.
  \item Determine which relationships are one-to-many and which are many-to-many.
  \item Propose intermediate tables if necessary.
\end{enumerate}
\end{frame}

%------------------------------------------------
\begin{frame}{Exercise 3: Employees and Dependents}
\small
A company wants to store information about employees and their dependents.

\vspace{2mm}
\begin{enumerate}
\smallitem
  \item Identify the strong entity and weak entity.
  \item Explain why one entity is a weak entity.
  \item Identify the identifying relationship.
  \item Suggest how to convert the model into relational tables.
\end{enumerate}
\end{frame}

%------------------------------------------------
\begin{frame}{Exercise 4: Hospital System}
\small
A hospital needs to manage patients, doctors, departments, surgeries, and medicines.

\vspace{2mm}
\begin{enumerate}
\smallitem
  \item Identify at least five entities.
  \item Identify at least five relationships.
  \item Classify some relationships by degree and cardinality.
  \item Indicate which relationships may require intermediate tables.
\end{enumerate}
\end{frame}

%================================================
\section{Summary and Keywords}

%------------------------------------------------
\begin{frame}{Lesson Summary}
\small
\begin{itemize}
\smallitem
  \item The ER Model is a conceptual model used for database design.
  \item An ER Diagram is the graphical representation of the ER Model.
  \item The three main components of the ER Model are entity, attribute, and relationship.
  \item An entity is an object or concept whose data must be stored.
  \item A strong entity has its own identifier; a weak entity depends on a strong entity.
  \item Attributes describe entities and include simple, key, composite, multivalued, and derived attributes.
\end{itemize}
\end{frame}

%------------------------------------------------
\begin{frame}{Lesson Summary (cont.)}
\small
\begin{itemize}
\smallitem
  \item Relationships represent connections between entities.
  \item Degree indicates the number of entity sets participating in a relationship.
  \item Cardinality indicates the maximum number of times an entity can participate in a relationship.
  \item Participation constraints indicate whether participation is mandatory or optional.
  \item The ERD drawing process includes identifying entities, relationships, attributes, primary keys, removing redundancies, and reviewing clarity.
\end{itemize}
\end{frame}

%------------------------------------------------
\begin{frame}{Key Terms}
\scriptsize
\begin{columns}[T]
\begin{column}{0.33\textwidth}
\begin{itemize}
\smallitem
  \item Entity-Relationship Model
  \item ER Model
  \item Entity-Relationship Diagram
  \item ERD
  \item Entity
  \item Entity Type
  \item Entity Set
  \item Entity Instance
  \item Strong Entity
  \item Weak Entity
\end{itemize}
\end{column}
\begin{column}{0.33\textwidth}
\begin{itemize}
\smallitem
  \item Attribute
  \item Simple Attribute
  \item Key Attribute
  \item Composite Attribute
  \item Multivalued Attribute
  \item Derived Attribute
  \item Relationship
  \item Relationship Type
  \item Relationship Set
  \item Degree
\end{itemize}
\end{column}
\begin{column}{0.33\textwidth}
\begin{itemize}
\smallitem
  \item Unary Relationship
  \item Binary Relationship
  \item Ternary Relationship
  \item N-ary Relationship
  \item Cardinality
  \item One-to-One
  \item One-to-Many
  \item Many-to-Many
  \item Participation Constraint
  \item Total/Partial Participation
\end{itemize}
\end{column}
\end{columns}
\end{frame}

%------------------------------------------------
\begin{frame}
\centering
\Huge Thank you!\\[0.5cm]
\Large Questions and Discussion
\end{frame}

\end{document}
